\section{Pipeline de entrenamiento iterativo}
\subsection{Desglose de la estructura de Self-Rewarding}
El circuito cerrado de autoevaluación de recompensa incluye tres roles: el Task Generator se encarga de derivar tareas a partir de las semillas de Self-Instruct, el Solution Generator da la respuesta inicial con Chain-of-Thought y el Reward Model ejecuta la verificación cruzada y calcula la recompensa. Weston considera que las tareas producidas con Self-Instruct se acercan más a las necesidades reales (código, matemáticas, evaluación), por lo que este paso no debe limitarse a repetir el conjunto de datos de instrucciones ya existente.

\begin{figure}[H]
\centering
\includegraphics[width=0.8\textwidth]{frame-01-selfreward.jpg}
\caption{El presentador muestra en un capítulo anterior el diagrama de conceptos de self-improving, que ilustra el pipeline de tres roles.\protect\footnotemark}
\end{figure}
\footnotetext{Intervalo de tiempo en el video: 00:02:35--00:02:37.}

\subsection{Self-Instruct y validación cruzada}
Para mantener la amplitud de las tareas, el pipeline usa Self-Instruct para generar por lotes nuevos prompts. Después de cada ronda de respuestas, el modelo genera una verified response que se empareja con el borrador inicial. Si el meta judge determina que ambas difieren, se usa la verified response para actualizar el modelo de recompensa. Este proceso se asemeja a la validación conversacional del ``Self-Feeding Chatbot'', pero se enfoca más en la cadena de razonamiento que en la simple respuesta conversacional.

\begin{importantbox}{La sinergia entre Self-Instruct y la verified response}
\begin{itemize}
    \item Self-Instruct: a partir de un número reducido de seed prompts, duplica automáticamente tareas de matemáticas, código y evaluación.
    \item Verified response: el modelo determina si se equivocó; aunque el borrador original ya contenga la respuesta escrita, de todos modos debe generar una versión nueva y verificada.
\end{itemize}
\end{importantbox}

\begin{figure}[H]
\centering
\includegraphics[width=0.78\textwidth]{frame-02-iterative.jpg}
\caption{Ilustración del pipeline de Self-Rewarding, con énfasis en la actualización iterativa de la verified response.\protect\footnotemark}
\end{figure}
\footnotetext{Intervalo de tiempo en el video: 00:36:00--00:36:02.}

\begin{figure}[H]
\centering
\includegraphics[width=0.75\textwidth]{frame-03-selfinstr.jpg}
\caption{Flujo de generación de tareas y validación de respuestas dentro de Self-Instruct.\protect\footnotemark}
\end{figure}
\footnotetext{Intervalo de tiempo en el video: 00:44:40--00:44:42.}

\subsection{Datos de entrenamiento y estrategia de fine-tuning}
Weston menciona repetidamente en la clase el ``Self-Feeding Chatbot'', los ``SelfT evaluators'' y las ``verified response iterations'', componentes que forman el eje principal de los datos de entrenamiento. En cada ronda de self-rewarding se agrega una generación alr (instruction, question, answer, reasoning) y, después, bajo la indicación del meta judge, se filtran las respuestas de mayor calidad. Vale la pena notar que, en los experimentos reales, la coherencia entre el modelo entrenado con tres rondas de self-rewarding y la verified response mejoró considerablemente, tal como dijo Weston: ``\textit{I presently love the simplicity of it that we didn't have to do this large amount of Chain of Thought where it sort of overthinks itself}''.

\begin{knowledgebox}{Estrategia de muestreo de Self-Instruct}
Self-Instruct parte de un número reducido de seeds hechas a mano y genera, de forma cuasi supervisada, tripletas de instruction → input → output. Weston recomienda cubrir en cada ronda cuatro tipos de tareas: lenguaje, matemáticas, programación y evaluación lógica, para evitar que el conjunto de datos de entrenamiento se concentre en un solo tipo de razonamiento.
\end{knowledgebox}

\subsection{Resumen de la sección}
El pipeline de entrenamiento de Self-Rewarding conecta al Task Generator, al Solution Generator y al Reward Model en un solo circuito cerrado. Self-Instruct garantiza la diversidad de las tareas, la Verified Response hace que el modelo de recompensa se concentre en los patrones de error, y Self-Feeding/el meta judge verifica además si el resultado vale la pena adoptarse.
